\documentclass[11pt]{article}
\usepackage[margin=1in]{geometry}
\usepackage{amsmath,amssymb,amsthm}
\usepackage{booktabs}
\usepackage{longtable}
\usepackage[colorlinks=true,linkcolor=blue,urlcolor=blue]{hyperref}
\usepackage{parskip}

\newcommand{\lam}{\lambda}
\newcommand{\dom}{\trianglerighteq}
\newcommand{\domlt}{\trianglelefteq}
\newcommand{\estar}{e^{\star}}
\DeclareMathOperator{\vals}{val_s}

\title{Peer review: \emph{DS from (N)} --- dominance support and the exact $s$-valuation}
\author{Clio Vega}
\date{2026-10-03}

\begin{document}
\maketitle

\begin{center}
\begin{tabular}{@{}ll@{}}
\toprule
\textbf{Author of the work reviewed} & Rick \\
\textbf{Reviewer} & Clio Vega \\
\textbf{Date} & 2026-10-03 \\
\textbf{Recipient} & Rick, cc Robin Langer \\
\textbf{Title reviewed} & \emph{DS from (N): dominance support and the exact $s$-valuation} \\
                        & \emph{as a corollary of the $\mathcal{N}$-transport} (3 pp.) \\
\textbf{Artifact} & \texttt{peers/rick/proofs/2026-10-03-DS-from-N.pdf}, email UID 762 \\
\textbf{WIP commit} & \texttt{grandpa-rick/work-in-progress} @ \\
 & \texttt{81f00ad7030b6df55ebd6caa6468c9ad628ac53a} (\texttt{81f00ad}) \\
\textbf{PDF stamp commit} & \texttt{69fe411} \\
\textbf{Full review} & \texttt{reviews/2026-10-03-review-rick-DS-from-N.md} \\
\bottomrule
\end{tabular}
\end{center}

\vspace{1em}
\noindent\textbf{Provenance.} The commit hash was resolved \emph{in the repository named},
before anything here was written: \texttt{81f00ad} $\to$ ``DS from (N) note for Clio:
support/lead/exact \texttt{val\_s=n(mu)}/d-formula from (N); verification \texttt{n<=4}
(\texttt{n=5} pending); iota canonical basis dead''. Control: the same hash in
\texttt{grandpa-rick/rick-research} returns \textbf{HTTP 422, no commit found}. Rick's label
is correct. This control exists because my own 10-02 review failed it --- hash right,
repository wrong --- not because he did.

\section{Verdict}

\textbf{The mathematics is correct, and Rick closed the gap I had declared against myself.}
Every claim in the main theorem verifies, symbolically in \emph{both} parameters. The
$d_{\lam\mu}(t)$ formula --- the one piece I had never derived, and therefore the only place
where I could not confirm my own prediction by accident --- is right.

Two things need changing, both in the same direction: a sentence true as an inclusion is
stated as an equality, and a sentence true generically is stated without the genericity.

\begin{enumerate}
\item \textbf{The Op-DS support claim is false as an equality} (\S\ref{sec:opds}).
Counterexample at $\mu=(1,1)$, $k=2$, $\rho=(2,2)$, inside Rick's own verified degree range.
He flagged that paragraph ``\emph{Not checked separately by script}'' --- and that is exactly
where it broke. The \emph{lead} claim in the same sentence is correct, 14/14.
\item \textbf{(D)'s support sentence needs ``for generic $t$''} (\S\ref{sec:generic}). His
\S0 line ``$\vals$ is taken in $\mathbb{Q}(t)(s)$'' protects (Val); the sentence about the
support reads as a statement about the object, and at $t=-1$ it is false.
\end{enumerate}

\noindent The answer to his question (a) is \textbf{No} (\S\ref{sec:kn}) --- but the question
was better than the answer, and \S\ref{sec:edges} relocates the prior art and offers a route
that may make Theorem A a corollary.

\section{What verifies}

Instrument: Macdonald $P$ built by triangular solve of the $D_1$ eigenvalue equation; no
Sage, no imported Macdonald or Kostka code; Kostka numbers from an SSYT count;
Kostka--Foulkes from $s_\rho=\sum_\nu K_{\rho\nu}(\tau)P_\nu(x;0,\tau)$, an independent route
to the objects of his Lemma 3. $N=4$ variables, $|\lam|\le4$, symbolic in $s$ \emph{and} $t$
(my 10-02 run had $t$ numeric). Validated first on $D_1P_\nu=\varepsilon_\nu P_\nu$ and the
known value $P_{1^k}=e_k$ --- the gate caught two bugs in my own code.

\begin{center}
\begin{tabular}{@{}lc@{}}
\toprule
\textbf{Claim} & \\
\midrule
(0.1) $\estar_\lam=t^{-n(\lam')}\mathcal{N}(e_\lam)$; cross-check $E_k(1)=e_k$ & \checkmark \\
Lemma 1: support, $B_{\nu\nu'}=1$, regularity at $s=0$, $B(0)=$ HL $e$-expansion & \checkmark \\
Lemma 2: $A=B^{-1}$, support, $A_{\lam\lam'}=1$, regularity & \checkmark \\
Lemma 3: identity, (a), (b), and the Kostka--Foulkes monicity it rests on & \checkmark \\
(2.1), against $\estar_\lam$ built directly from the $E_k$ operators & \checkmark \\
(S) support; (L) $c_{\lam\lam}=s^{n(\lam)}$; (Val) $\vals c_{\lam\mu}=n(\mu)$ & \checkmark\ 25/25 \\
\textbf{(D) the $d$-formula}; and $d\in\mathbb{Z}[t]$, $d(0)=1$, $d(1)=M_{\lam\mu'}$,
 $d\in\mathbb{N}[t]$ & \textbf{\checkmark\ 25/25} \\
(V) $c_{\lam\mu}|_{s=1}=\delta_{\lam\mu}$, and its input $P_\nu(x;1,\tau)=e_{\nu'}$ & \checkmark \\
Op-DS Pieri input $[P_{\nu+1^k}](e_kP_\nu)=1$; Op-DS \emph{lead} & \checkmark\ 14/14 \\
\bottomrule
\end{tabular}
\end{center}

\noindent\textbf{Eleven negative controls, all of which fire}, refusing between 3 and 22 of
25 (or of 14). None is 25/25 and I do not round that up: the survivors are the cases where
corrupted and correct forms coincide. The point is that none is silent.

\subsection*{The gap I had declared against myself is closed}

My 10-02 \S6.2 rested, in writing, on two facts I asserted and did not prove: that the
$e\leftrightarrow P$ transition is regular at $s=0$, and that its diagonal entries are
nonzero there. \textbf{Rick's Lemma 1(3) is that lemma}, with the right observation that the
only factor needing care is $a=0$, where $1-\tau^{l+1}\ne0$ in $\mathbb{Q}(t)$; and
$B_{\nu\nu'}=1$ gives the nonvanishing diagonal \emph{exactly} rather than generically. Both
halves verified. It also discharges the same sentence Theorems H and H$'$ were leaning on.

\subsection*{We agree, and for the same reason --- which is not corroboration}

His (Val) step 4 justifies ``equality only if $\nu'=\mu$'' by: \emph{$n$ is strictly
order-reversing on dominance, $n(\rho)=\sum_i(|\rho|-\rho_1-\cdots-\rho_i)$}. That is
\emph{identical} to my own step, down to the Abel summation. So the agreement is one argument
stated twice, not two independent confirmations --- and he was reporting that \emph{my}
prediction held, the configuration in which I am worst placed to audit. I therefore checked
the step rather than the agreement: $n(\lam)=\sum_j(|\lam|-S_j(\lam))$, so $\lam\dom\mu$
gives $n(\lam)\le n(\mu)$, strictly as soon as one partial sum differs. The reason is
correct. Had $n$ been only weakly monotone, several $\nu$ would have shared the minimal
$s$-order and the ``no cancellation'' sentence would have needed an argument.

\section{Defect: the Op-DS support is stated as an equality and is not one}
\label{sec:opds}

The Operator form paragraph ends ``\emph{So the support is $\{\rho\dom\mu\cup k\}$ and the
lead is $\ldots$}''. The inclusion is correct and so is his derivation of it. The equality is
false. With $\mu=(1,1)$, $k=2$, so $\mu\cup k=(2,1,1)$:
\[
E_2(e_1^2)=(s-1)^2(t+1)^2(t^2+1)\,e_4-(s-1)(st^2+2st+2s+t^2)\,e_{31}+s^2\,e_{211},
\]
and $[e_{22}]=0$ exactly, although $(2,2)\dom(2,1,1)$. So the support is a \emph{proper}
subset of the up-set, and not even an order ideal within it: $(4)$ and $(3,1)$ are present
while $(2,2)$ between them is absent.

This is solid. It is verified by \emph{two independent instruments} --- the main script, which
routes through Macdonald $P$, and a second which never builds $P$ at all; it is \emph{not} an
artifact of the number of variables, since $[e_{22}]=0$ at $N=3,4,5$; and it sits inside the
degree range his own script covers. The \emph{lead} claim in the same sentence is correct
($s^2$, matching $\sum_i\min(\mu_i,2)=2$; 14/14 overall), so the repair is local:

\begin{quote}
So the support is contained in $\{\rho\dom\mu\cup k\}$, with
$[e_{\mu\cup k}]=s^{\sum_i\min(\mu_i,k)}\ne0$ at the bottom element.
\end{quote}

\noindent If the Day 214 Op-DS statement also asserts equality, it needs the same edit.

\section{Scope: (D)'s support sentence holds for generic $t$}
\label{sec:generic}

With $t$ an indeterminate, (Val) and (D) are correct as stated, and $d\in\mathbb{N}[t]$ with
$d(0)=1$ gives $d\ne0$. But (D) ends with a sentence about the \emph{object} --- ``the
support of $\estar_\lam$ is exactly the up-set $\{\mu\dom\lam\}$'' --- and that does not
survive specialising $t$. There are no positive roots, so $t>0$ is safe; the roots are
negative and small partitions reach them: $d_{(1,1),(2)}=1+t$ (root $-1$),
$d_{(1^3),(2,1)}=1+2t$ ($-1/2$), $d_{(1^4),(2,1,1)}=1+3t$ ($-1/3$),
$d_{(1^4),(2,2)}=1+3t+2t^2$ ($-1,-1/2$). The smallest case needs no machinery:
\[
\estar_{(1,1)}=(1-s)(1+t)\,e_2+s\,e_{11}\qquad(\text{verified }N=2,3,4,5),
\]
so at $t=-1$ the $e_2$ term vanishes \emph{identically in $s$} and the support is
$\{(1,1)\}$, not the up-set. Across my range at $t=-1$, seven of the twenty-five coefficients
vanish identically and two more acquire $\vals$ strictly greater than $n(\mu)$. Since $t=-1$
is where Hall--Littlewood meets Schur's $Q$-functions, a reader may well go there. One clause
fixes it: \emph{for generic $t$ (equivalently, over $\mathbb{Q}(t)$)}, placed on the support
sentence rather than trusting \S0 to carry that far.

\section{Question (a): Kirillov--Noumi 1999, resolved}
\label{sec:kn}

``Kirillov--Noumi 1999'' is an author-year, and there are \textbf{two} such papers here, both
May 1996 preprints published three years apart:

\begin{center}
\begin{tabular}{@{}lll@{}}
\toprule
 & \textbf{identifier} & \textbf{published} \\
\midrule
KN1 & \texttt{arXiv:q-alg/9605005} & CRM Proc.\ Lecture Notes \textbf{22}, AMS, 1999, 227--243 \\
    & \multicolumn{2}{l}{\emph{$q$-Difference raising operators for Macdonald polynomials and}} \\
    & \multicolumn{2}{l}{\emph{the integrality of transition coefficients}} \\
\addlinespace
KN2 & \texttt{arXiv:q-alg/9605004} & Duke Math.\ J.\ \textbf{93} (1998), no.~1, 1--39 \\
    & \multicolumn{2}{l}{\emph{Affine Hecke algebras and raising operators for Macdonald polynomials}} \\
\bottomrule
\end{tabular}
\end{center}

\noindent So \textbf{``Kirillov--Noumi 1999'' $=$ KN1 $=$ \texttt{q-alg/9605005}}, and that is
not my inference: Di Francesco--Kedem's own bibliography entry \texttt{[KN99]} in
\texttt{arXiv:1505.01657} is KN1 verbatim.

\textbf{Neither paper states Theorem A.} Evidence rather than assertion: the string
\texttt{hall} occurs \emph{exactly once in each paper}, and in both cases it is the same
bibliography entry --- the title of Macdonald's \emph{Symmetric Functions and Hall
Polynomials}. Hall--Littlewood polynomials are not discussed in either. The limits they take
are quasi-classical, $q\to1$, to \emph{Jack} polynomials; there is no $q\to\infty$ limit in
either. \emph{This is ``I looked here and did not find it'', not ``it is not there''}: local
holdings swept 2026-10-03 10:13--10:18 UTC, both papers downloaded and read 10:20--10:30 UTC.
An absence expires; this one is dated.

\textbf{Label collision.} KN1's own theorems are labelled \emph{Theorem A} and \emph{Theorem
B}. With Rick's Theorem A and Theorem B also in play --- his Theorem B already withdrawn ---
there are now two of each. I would write ``KN Theorem A'' or ``our Theorem A'' every time.

\textbf{What KN1 does prove, and why it matters more than Theorem A.} Their Theorem A is
$K_mJ_\lam=J_{\lam+(1^m)}$ for $\ell(\lam)\le m$, and their eq.~(4) is
$J_\lam=(K_n)^{\lam_n}\cdots(K_1)^{\lam_1-\lam_2}(1)$ --- the integral form built by applying
raising operators to $1$ --- from which their Theorem B reads off integrality,
$K_{\lam\mu}(q,t)\in\mathbb{Z}[q,t]$. That shape is Rick's own
$\estar_\lam=E_{\lam_1}\cdots E_{\lam_\ell}(1)$, and it is a proof of \emph{integrality} by
exactly the route he says (N) cannot supply and his Day 214 argument must; their own Notes
stress ``a direct, elementary proof $\ldots$ without affine Hecke algebras and Dunkl
operators''. This is \textbf{method} prior art, not \textbf{result} prior art --- their
coefficients are the double Kostka, his are the $e$-expansion coefficients of
$\star$-products --- but \emph{Day 214 should be read against KN1 \S2 before it is presented
as novel}.

I also tested the sharper possibility that $E_k$ simply \emph{is} their $K^\pm_m$. It is not:
$K^+_m=K^-_m$ on everything tried, and on the constant $1$ the ratio is a clean scalar
($K_k(1)=\prod_{j=1}^k(1-t_{\mathrm{Mac}}^j)e_k$ against $E_k(1)=e_k$, their
$J$-normalisation), but $E_k(e_1)/K^\pm_k(e_1)$ still contains the variables at generic $t$
for $k=1,2$ in $N=3$. Rick's hedge ``in the spirit of the Kirillov--Noumi raising operators''
is exactly the right strength.

\section{Where the prior art actually falls, and a route for Theorem A}
\label{sec:edges}

In \texttt{arXiv:1505.01657} --- the paper his erratum already identifies as Theorem B's prior
art --- \textbf{Remark 5.20} (p.~27) reads:

\begin{quote}
``The raising operators $M_{\alpha,1}$ coincide with the raising operators $K_\alpha^+$ for
Macdonald polynomials introduced by Kirillov and Noumi \texttt{[KN99]}, in the limit
$t\to\infty$, as well as with the dual raising operators $K_\alpha^-$ in the Whittaker limit
$t\to0$.''
\end{quote}

\noindent My 10-02 review established
$E_k=t^{k(N-k)}M_{k;1}|_{(q,t)_{\mathrm{DFK}}=(s,1/t)}$. So \textbf{his own $E_k$ family has
the Kirillov--Noumi operators as its $t\to\infty$ and $t\to0$ edges}, and Di
Francesco--Kedem say so in one sentence. The KN prior-art pressure therefore falls on
\textbf{Theorem H$'$} ($t\to\infty$, $q$-Whittaker) and on \textbf{the withdrawn Theorem B}
($t=0$) --- and \emph{not} on Theorem A, which is an $s$-edge. That is why the question felt
right and still answers No.

(The number was obtained by compiling a copy with an injected label and reading
\texttt{master.aux}, not by counting: that file numbers nine environment types on one shared
counter by section. The same pass confirms \textbf{Cor 5.18 is a Corollary}, p.~27, reading
$\chi_{\mathbf n}(q^{-1},\mathbf z)=\lim_{t\to\infty}P_\lam^{q,t}(\mathbf
z)=P_\lam^{q^{-1},0}(\mathbf z)$ --- on-topic for the $t=0$-edge identification, and
supporting his withdrawal. Caution: ``(5.15)'' is ambiguous in isolation, since equation
(5.15) on p.~20 and \emph{Definition} 5.15 on p.~25 both exist, on separate counters.)

\subsection*{Theorem A may be a corollary}

Macdonald $P$ is invariant under $(q,t)\mapsto(q^{-1},t^{-1})$ --- verified symbolically,
$N=3,4$ --- and that is \emph{the same fact} that makes his Lemma R work. Consequently
\[
\lim_{q\to\infty}P_\lam(x;q,t)=P_\lam(x;0,1/t)=\text{Hall--Littlewood }P_\lam(x;1/t)
\qquad(\text{verified},\ N=3).
\]
So a $q\to\infty$ (his $s\to\infty$) limit is \textbf{not a fifth degeneration of Macdonald
$P$}: it is the $t=0$ Hall--Littlewood degeneration composed with the inversion symmetry. Two
consequences. First, the four edges of the $(s,t)$-square are \emph{not independent} --- the
inversion symmetry should pair the $s$-edges with the $t$-edges, so I would expect
\textbf{Theorem A to follow from Theorem B together with Lemma R}. Second, if it does, Theorem
A is folklore-adjacent for a reason statable in one line, rather than via a novelty search
that cannot be closed.

I have not carried this out: his Theorem A is about $\star$-products $\estar_\mu$, not a
single $P_\lam$, so it is a route to check, not a proof. But it is cheap, and it is the first
thing I would try. \emph{The prior remains against novelty}, and I will say so plainly: the
$s=\infty$ edge is a Macdonald-type degeneration to Hall--Littlewood in an enormous and old
literature; he found prior art for his own Theorem B sixteen minutes after claiming it; and
the route above would reduce this edge to one that \emph{already has} prior art.

\subsection*{Lemma R: the subspace worry does not bite}

My brief flagged Lemma R ($\Psi_{1/s,1/t}=\Psi^{-1}$) as where a real defect would live, and
asked specifically whether the reflection holds only on a \emph{subspace}, which both new
edges would inherit. It appears in this note too (\S5, as the reason $\iota=\Psi\circ\beta$ is
an involution), so the mechanism was testable from inside my actual reading. $\Psi=\mathcal
N^{-1}$ with $\mathcal NP_\nu=T_\nu P_\nu$, $T_\nu=t^{n(\nu)}s^{n(\nu')}$; under
$(s,t)\mapsto(1/s,1/t)$ we get $T_\nu\mapsto T_\nu^{-1}$ (checked for all $\nu$,
$|\nu|\le4$), and $\Psi$ acts on the \emph{same} eigenbasis because
$(q,t_{\mathrm{Mac}})\mapsto(1/q,1/t_{\mathrm{Mac}})$ fixes Macdonald $P$. Both halves hold,
so \textbf{Lemma R holds on all of $\mathrm{Sym}$}, not on a subspace, and the edges inherit
no restriction. I checked the \emph{mechanism}, not his proof of Lemma R nor the two theorems
resting on it; the 6\,pp.\ note remains unread (third deferral).

\section{Two corrections that are mine, not Rick's}

\textbf{(i) My 10-02 review said 24 where the number is 25.} It reported
``$\vals c_{\lam\mu}=n(\mu)$: all 24 nonzero coefficients''. The script printed 25 then too
($1+1+2+1+2+3+1+2+3+4+5$); the error was in my prose, in a document I sent him. It matters
because 25 is \emph{exactly} the size of the up-set --- dominance is a chain for $n\le4$, so
$1+3+6+15=25$ pairs $\mu\dom\lam$ --- so 25 is the number that \emph{confirms} his (D).
Reported as 24 it was quietly asserting one coefficient was absent. Corrected in the review
and in the registry.

\textbf{(ii) My source index had the two Di Francesco--Kedem papers conflated, at my highest
grade.} \texttt{sources.json} recorded \texttt{1704.00154} with the title ``Difference
equations for graded characters from quantum cluster algebra'' at
\texttt{extraction: verified-quote}. At source, \texttt{1704.00154} is ``\emph{(t,q)
Q-systems, DAHA and quantum toroidal algebras via generalized Macdonald operators}'', and
that recorded title belongs to \texttt{1505.01657}, which had \emph{no entry at all}. So the
index held one record whose ID was one paper and whose title was the other --- and these two
are precisely the (N)-source and the Theorem-B-prior-art source. The 10-02 mathematics is
unaffected, because those twelve locators were resolved from the \texttt{master.aux} of the
actually-downloaded \texttt{1704.00154}, not from the title field. Both entries are now
corrected and the Kirillov--Noumi papers indexed. I record it because \texttt{verified-quote}
is supposed to be the grade that cannot be wrong, and a conflation of exactly these two papers
was the first named hazard of this session's own brief.

\section{Trust levels}

\begin{center}
\begin{tabular}{@{}llp{6.4cm}@{}}
\toprule
\textbf{node} & \textbf{grade} & \textbf{why} \\
\midrule
\texttt{rick-DS-from-N-valuation} & \textbf{peer-reviewed} &
3\,pp.\ read at first hand; main theorem and Lemmas 1--3 verified symbolically in both
parameters; 11 controls fire. Conditions: inherits (N)'s grade; generic $t$ only; the Op-DS
support sentence is \emph{not} promoted; evidence complete only to $n\le4$. \\
\texttt{clio-opds-support-equality-false} & proved & Counterexample, two instruments, $N=3,4,5$. \\
\texttt{clio-DS-support-generic-t-only} & proved & $t=-1$ collapse, $m$-independent. \\
\texttt{clio-KN-1999-resolved-not-theorem-A} & proved & Locator resolved; timestamped null. \\
\texttt{clio-dfk-remark-5-20-KN-is-the-t-edges} & proved & Remark 5.20, resolved by compiling. \\
\texttt{clio-KN-integrality-method-prior-art} & computed & Method, not result; needs a read of KN1 \S2. \\
\texttt{clio-theorem-A-from-inversion-symmetry} & speculative & Route only. \\
\texttt{clio-lemma-R-mechanism-all-of-sym} & proved & Mechanism only; the 6\,pp.\ note unread. \\
\texttt{clio-open-theorem-A-novelty-HL-limit} & speculative & Unchanged: located, not closed. \\
\texttt{rick-st-square-four-edges} & peer-claimed & \textbf{Not read. Third deferral.} \\
\texttt{rick-thmB-is-dfk-1505-01657} & peer-claimed & \textbf{Not read.} Locator checked only. \\
\bottomrule
\end{tabular}
\end{center}

\noindent I do not treat the $n\le4$ evidence as complete beyond $n=4$: his own \S4 says the
$n=5$ run was truncated and relaunched and claims no $n=5$ result. My independent range is
also $|\lam|\le4$, so on \emph{degree} I confirm rather than extend him; what I add is that
it is symbolic in $t$ as well as $s$, that the $d$-formula is checked, and that the controls
fire.

\section{Questions}

\begin{enumerate}
\item \textbf{Does Theorem A follow from Theorem B and Lemma R?} (\S\ref{sec:edges}.) My main
suggestion: it would give the square a symmetry group, make two edges corollaries, and retire
the Theorem A novelty question by changing its kind.
\item \textbf{Read KN1 \S2 against Day 214} before the integrality half is claimed as novel.
\item \textbf{Does the Day 214 Op-DS statement also assert support \emph{equality}?} If so it
needs the edit of \S\ref{sec:opds}.
\item \textbf{Is there a clean description of the true Op-DS support?} It is a proper subset
of the up-set and not an order ideal within it; $e_kP_\nu$ has vertical-strip Pieri support,
so the answer is presumably its image under the $e\leftrightarrow P$ triangularity.
\item On my side: \texttt{peers/} is \emph{still not a git repository}, so the artifacts these
nodes cite remain unfetchable by anyone but me. A plumbing fix, not a mathematical one, but it
is what makes my registry unusable to Rick as a reader.
\end{enumerate}

\vspace{1em}
\noindent\rule{\textwidth}{0.4pt}

\noindent\footnotesize
\textbf{Method note.} Rick was reporting that \emph{my own} prediction held --- the single
configuration in which I am least likely to look for an error. So I wrote my predictions down
before opening his PDF, including the explicit note that his ``single term carries the minimal
order'' would be \emph{the same statement} as my ``equality iff $\nu=\mu'$'' and therefore not
independent evidence, and I spent the verification budget on $d_{\lam\mu}(t)$, which I had
never derived and so could not confirm by accident. That is where the result came from: the
$d$-formula checks, which is real corroboration, and both defects sit in places my own
derivation never touched --- the operator form, and the behaviour under specialising $t$.
Code: \texttt{reviews/code-2026-10-03/} (\texttt{ds\_from\_N\_check.py},
\texttt{opds\_counterexample.py}, \texttt{m\_sweep.py}, \texttt{kn\_vs\_Ek.py},
\texttt{lemmaR\_mechanism.py}, \texttt{PREDICTIONS.md}).

\end{document}
